\chapter{Convergence}

The morning bicycle ride felt like riding into an abyss that began ten feet ahead. Fog spread all around, and the wind was the worst of it. The road along the dry riverbed, which had once carried the summer loo, now lay hidden under a white veil. Shiva usually reached school half an hour early, as did most of his friends. Sitting beside a fire before the morning assembly had become their routine. The hostel students joined them too, sometimes, once they had bathed in the chilling water.

``Why didn't you come for the volleyball games yesterday?'' Vijay asked, slowly pushing a piece of broken bench-wood into the fire.
``No, it's just that I've joined tuition with Mr. Verma after school,'' Shiva replied, his eyes on the flames as they gently grew.
``Oh—so that's why you weren't at Robin sir's coaching today. But wasn't Robin sir teaching you well enough already? He liked you so much, too.''
``I know. I told him yesterday, when I was there, that I was leaving,'' Shiva said, stretching his hand forward to push another log into the fire. As he did, the sleeve of his blazer slipped back and his arm caught the light. The letters had faded almost entirely now; only a darker circle remained. He pulled the sleeve down quickly.
``I know why you joined there,'' Rajveer said with a mischievous grin. ``Just take care of it cleanly this time. Make sure it doesn't end up like the last.''

``Don't worry,'' Shiva said with a smile. A bell rang out, echoing across the campus, and they stood up for the assembly. Gauri led the prayer every morning; she was the school's lead singer. After all the gaps, his mornings had once again begun with her voice.

Shiva stayed back in class until everyone had left. Only once it was empty did he pick up his bag and walk up to the second floor. His friends were still on the ground below, playing. He glanced at them as he crossed the corridor, and Rajveer, happening to look up just then, caught his eye and gave him a thumbs up. Shiva smiled and walked on. Entering the room, he found a closed tiffin on the desk, Roshini and Gauri seated on either side of it, both staring at him. He paused. ``What happened?'' he asked.
``She won't let me eat,'' Gauri said, drama in her voice.
``It's just that she's made gajar-halwa for the first time. I wanted you to taste it too, not only me. Suppose it's bad—really bad, salt-instead-of-sugar bad. How can I tease her about it all on my own? I need a backup, right?'' Roshini said, determined.
``Can I join, then?'' Shiva asked Gauri, pointing to the tiffin.
``Sure. But if you two overdo it, remember—I'll make halwa out of you as well,'' Gauri replied.

Shiva took the spoon resting on top of the tiffin, and Roshini opened it. He scooped out a little and put it in his mouth. His face went pale and still; he even stopped chewing. The girls' heartbeats stopped with him. ``Is there really salt in it?'' Roshini exclaimed. Shiva nodded. Roshini took the spoon and tried a taste of her own. She froze too. Now the pressure fell on Gauri. ``What? No...'' she shot back. ``I tasted it this morning—it wasn't this bad. What's wrong?''

Shiva and Roshini broke into laughter, and it pulled Gauri into a smile too. ``No, it's genuinely good. I wish you'd bring it more often,'' Shiva said, calming her. ``Oh, thanks. Let's see how many more days we have carrots at home,'' she said, smiling along with them.

They finished the box just before Mr. Verma arrived, then took the seats facing him—Gauri on the left, Roshini in the middle, and Shiva on the right.

Days passed slowly. Mr. Verma arrived for tuition about half an hour after school let out, and the three of them made the most of that gap—sharing food, usually Gauri's tiffin, then heading down to the ground for a game of badminton, or letting Shiva teach them how to handle a volleyball. His friends had slowly dropped the habit of playing after school as the winter turned harsher, so the court lay open now for the three of them, for the few minutes they had.

One such winter evening, Shiva returned home from school. Nani sat near the small fire, which had almost faded. He stood his cycle in the backyard veranda, and after changing his clothes he came back and settled down beside her. ``The fire's almost gone. Should I add some wood?'' he asked. ``Yes. But you'll have to split the log first—it's too large,'' Nani replied, slowly turning over the remaining embers with a small stick. Shiva stood up, fetched a smaller log and a hand axe, and started splitting it.

``Mama has brought everything for tomorrow's pooja. Before it gets dark, go to Pandit ji's house and remind him about tomorrow,'' Nani said.
``Will Pandit mama come, or his son?'' Shiva asked, gathering the smaller pieces he'd split and setting them to one side.
``His son will come. He's grown old now, and much weaker. And he's been unwell with the cold,'' Nani replied.
``Okay, then I'll just call him—I have his number in my phone,'' Shiva said promptly. ``No, go to him in person. And meet his father too; ask after his health. He used to come to our home so often,'' Nani replied, turning the embers again. ``When you were little, he'd play with you for hours. You were always asking him to play horse-horse, and he would carry you on his back,'' Nani continued. ``But his health has gotten the better of him now. Otherwise he was always so active—from his priestly duties to the fields.''
``All right, Nani. I'll go to his house,'' Shiva said, gathering a few more pieces and laying them in the dying fire.

Soon the fire brightened. Shiva stood up and set off for Pandit ji's house. He reached it before long and found the veranda empty. He called his friend Vinay, Pandit ji's son, and asked him to come out. Vinay came around from the next street, and they shook hands. ``What happened, Shiva?'' Vinay asked. ``Oh—I came to remind you about tomorrow's pooja,'' he replied. ``Yes, don't worry, I remember. I'll come in good time,'' Vinay said with a smile. ``You could have just called me,'' Vinay added. ``No, I also wanted to see your father—I heard he's unwell,'' Shiva said. ``Yeah, his knee's swollen from the cold,'' Vinay replied. ``He had an operation on it once, and now it aches whenever the weather turns. And this time he's caught a cold as well.'' ``Oh. Have you taken him to a doctor?'' Shiva asked. ``Yes. The doctor's given medicines. But mostly it's about looking after him.''

``Where is he?'' Shiva asked.

Vinay led the way inside, and they crossed the veranda into a gallery. It was dark, with almost no light except what came in from either end, and it opened onto an inner courtyard that lay open to the sky. They turned left, and Vinay slid a door open. The room was lit by a single bulb, hanging from a holder tied to a nail on the wall. Pandit ji lay on the bed under a blanket, his wife seated beside him with a polythene bag of medicines in her lap. ``Namaste, Mami,'' Shiva said, joining his hands. ``Be happy, son,'' she replied. ``How are you? It's been so long since I saw you,'' she added, setting the polythene bag aside near the pillow. ``I'm good, Mami. How are you? And how is Mama?'' he asked. ``He's better now. I was just giving him his medicines—the ones before dinner,'' she said, her voice growing heavy, almost breaking into tears. ``Look who's come,'' she said, gently tapping the blanket. A face slowly emerged from beneath it, rising against the headrest of the bed. ``Oh—Shiva? How are you?'' he asked, joy in his voice. ``I'm good, Mama.'' Vinay pulled a chair over from the side and sat down. ``How are your Nani and Nana?'' Pandit ji asked. ``They're well too.''
``I hope Rama ji makes you well again soon—so I don't have to rely on Vinay for the pooja,'' Shiva said, trying to lift the mood. Vinay smiled, and so did Mami. ``Definitely. But for now, you'll have to make do with just him,'' Pandit ji said, breaking into a light laugh. ``What other choice do I have?'' Shiva added, with a glance toward Vinay.